\documentclass[11pt,a4paper]{article}

\usepackage[utf8]{inputenc}
\usepackage[T1]{fontenc}
\usepackage[ngerman]{babel}
\usepackage[a4paper,left=22mm,right=22mm,top=18mm,bottom=20mm]{geometry}
\usepackage{xcolor}
\usepackage{hyperref}

% Trumpf-inspirierte Farben und Schrift passend zum Lebenslauf.
\definecolor{CVPrimary}{HTML}{0033BA}
\definecolor{CVAccent}{HTML}{4D7CE6}

\hypersetup{colorlinks=true,urlcolor=CVPrimary}
\renewcommand{\familydefault}{\sfdefault}
\pagestyle{empty}
\setlength{\parindent}{0pt}
\setlength{\parskip}{0.65em}
\setlength{\emergencystretch}{2em}

\newcommand{\Briefdatum}{\today}

\begin{document}
\color{black}

{\raggedleft \Briefdatum\par}
\vspace{5mm}
{\large\bfseries\textcolor{CVPrimary}{Bewerbung als Trainee Technology \& Innovation (w/m/d)}}\par
{\small Referenznummer R00042851}\par
{\color{CVAccent}\rule{\linewidth}{0.6pt}}\par
\vspace{1mm}

Sehr geehrte Frau Wachter, 

nur wer neue Technologien mutig in die Fertigung bringt, bleibt in der Industrie vorn. Genau daran möchte ich im Trainee-Programm bei TRUMPF mitarbeiten.

Im dualen Bachelorstudium der Wirtschaftsinformatik bei Bosch durchlief ich verschiedene IT-nahe Abteilungen, unter anderem im R\&D-Umfeld zur Verbesserung von Parkassistenten, und absolvierte ein Praxissemester in Polen. Dabei lernte ich, mich zügig in neue Themen und Teams einzuarbeiten. Zudem präsentierte ich regelmäßig auch nicht-technische Themen anforderungsgerecht für fachfremde Zielgruppen.

Im Masterstudium Machine Learning \& Data Analytics an der Hochschule Aalen vertiefte ich meine Kenntnisse in Deep Learning und Computer Vision. Bei ZEISS SMT entwickelte ich als Werkstudentin einen Reinforcement-Learning-Agenten für industrielles Scheduling und analysierte die Maschinenauslastung. Aktuell arbeite ich als wissenschaftliche Mitarbeiterin an einem datengetriebenen Verfahren zur Restnutzungsdauer-Vorhersage von Bandsägen. So verbinde ich KI-Methoden mit konkreten Fragestellungen aus Maschinen und Produktion.

Bei Selfnet e.,V., einer studentischen Initiative zum Betrieb der Internetinfrastruktur in Studentenwohnheimen, wirke ich an der Entwicklung des Mitgliederverwaltungssystems mit und sammle praktische DevOps-Erfahrung. Als Mentor begleite ich neue Mitglieder, darunter viele internationale Studierende, beim Einstieg. Diese Teamarbeit und Eigeninitiative möchte ich bei Trumpf einbringen.

Bei InnoConsult e.\,V., einer studentischen Unternehmensberatung, führte ich Design-Thinking-Workshops durch und entwickelte Strategien für Startups. Bei Selfnet e.\,V., einer studentischen Initiative zum Betrieb der Internetinfrastruktur in Studentenwohnheimen, wirke ich an der Entwicklung des Mitgliederverwaltungssystems mit und begleite als Mentorin neue Mitglieder. Mich motiviert dabei vor allem, Ideen mit anderen zusammen von der Skizze bis zur funktionierenden Lösung zu bringen.

 Besonders reizt mich an dem ausgeschriebenen Programm, in wechselnden Einsätzen verschiedene Bereiche des Unternehmens kennenzulernen und früh Verantwortung für Technologie- und Innovationsthemen zu übernehmen.

Ich freue mich darauf, Sie in einem persönlichen Gespräch von meiner Motivation zu überzeugen und mehr über die Stationen des Programms zu erfahren.

Mit freundlichen Grüßen

\vspace{6mm}
Natalie Böhm

\end{document}


